%% kapa: exam=1 round=23 year=2012 period=2 form=A subject=accounting subject_ko=회계학 seq=17 number=17
%% source: Q-Net 감정평가사 기출문제 (cst00309) artlSeq=5052580 file=raw/1008/5052580/2012년 제23회 감정평가사 제1차시험 기출문제(2교시 A형).pdf
%% conversion: kordoc 4.24.51 -> tools/gen1.py (auto); review=proofread-sonnet; flags=none
\begin{question}{17}
충당부채, 우발부채 및 우발자산에 관한 설명으로 옳지 않은 것은?
\begin{choices}
  \item 소송이 진행중인 경우는 보고기간 말에 현재의무가 존재하지 아니할 가능성이 높더라도 경제적 효익이 내재된 자원의 유출가능성이 아주 낮지 않는 한 충당부채로 공시한다.
  \item 우발부채는 경제적 효익이 내재된 자원의 유출을 초래할 현재의무가 있는지의 여부가 아직 확인되지 아니한 잠재적 의무이다.
  \item 충당부채는 현재의무이고 이를 이행하기 위하여 경제적 효익이 내재된 자원이 유출될 가능성이 높고 당해 금액을 신뢰성 있게 추정할 수 있으므로 부채로 인식한다.
  \item 충당부채로 인식하는 금액은 현재의무를 보고기간 말에 이행하기 위하여 소요되는 지출에 대한 최선의 추정치이어야 한다.
  \item 미래에 전혀 실현되지 아니할 수도 있는 수익을 인식하는 결과를 초래할 수 있기 때문에 우발자산은 재무제표에 인식하지 아니한다.
\end{choices}
\end{question}
\answer{1}
